% ============================================================================
% SEC05.TEX - Section 2.5: Asset Store dan Package Management
% ============================================================================

\section{Unity Asset Store dan Package Management}

\begin{learningoutcomes}
\item Menggunakan Unity Asset Store untuk mencari assets
\item Memahami package management di Unity
\item Mampu menginstall dan mengelola packages
\item Memahami perbedaan Asset Store dan Package Manager
\end{learningoutcomes}

\subsection{Unity Asset Store}

Unity Asset Store adalah marketplace untuk:
\begin{itemize}
\item 3D models dan textures
\item 2D sprites dan UI elements
\item Audio assets
\item Scripts dan tools
\item Complete projects dan templates
\end{itemize}

\subsection{Menggunakan Asset Store}

\begin{enumerate}
\item Buka Window $\rightarrow$ Asset Store
\item Browse atau search assets
\item Pilih asset dan klik "Add to My Assets"
\item Download melalui Unity Hub atau Package Manager
\item Import ke project
\end{enumerate}

\subsection{Package Manager}

Package Manager mengelola:
\begin{itemize}
\item Unity packages (built-in)
\item Asset Store packages
\item Custom packages
\end{itemize}

\textbf{Akses Package Manager}: Window $\rightarrow$ Package Manager

\subsection{Types of Packages}

\begin{enumerate}
\item \textbf{Unity Registry}: Packages resmi dari Unity
\item \textbf{My Assets}: Assets dari Asset Store
\item \textbf{In Project}: Packages yang sudah terinstall
\item \textbf{Local}: Packages dari local disk
\end{enumerate}

\subsection{Installing Packages}

\begin{enumerate}
\item Buka Package Manager
\item Pilih source (Unity Registry, My Assets, dll)
\item Pilih package
\item Klik "Install"
\end{enumerate}

\subsection{Best Practices}

\begin{itemize}
\item \textbf{Review Before Install}: Baca deskripsi dan review
\item \textbf{Check Compatibility}: Pastikan kompatibel dengan versi Unity
\item \textbf{Organize}: Simpan assets di folder yang sesuai
\item \textbf{Update Regularly}: Update packages untuk bug fixes
\end{itemize}

\begin{catatan}
Asset Store sangat membantu untuk mempercepat development, tetapi pastikan untuk memahami bagaimana assets bekerja sebelum menggunakannya dalam project.
\end{catatan}

\subsection{Ringkasan Section}

\begin{itemize}
\item Asset Store menyediakan berbagai assets untuk mempercepat development
\item Package Manager mengelola packages dalam project
\item Ada berbagai jenis packages: Unity Registry, Asset Store, Local
\item Penting untuk review dan organize assets dengan baik
\end{itemize}
